\documentclass[acmsmall,review]{acmart}
\usepackage{listings}
\usepackage{xcolor}
\usepackage{booktabs}
\usepackage{microtype}

% OCaml listings style
\lstdefinelanguage{OCaml}{
  keywords={let,in,fun,match,with,type,module,sig,struct,
            and,or,not,if,then,else,begin,end,val,mutable,
            open,include,functor,of,rec,when},
  keywordstyle=\color[rgb]{0.0,0.4,0.7}\bfseries,
  commentstyle=\color[rgb]{0.4,0.4,0.4}\itshape,
  stringstyle=\color[rgb]{0.2,0.6,0.2},
  morecomment=[s]{(*}{*)},
  morestring=[b]",
  sensitive=true
}
\lstset{
  language=OCaml,
  basicstyle=\ttfamily\small,
  frame=single,
  framerule=0.4pt,
  rulecolor=\color[rgb]{0.8,0.8,0.8},
  backgroundcolor=\color[rgb]{0.97,0.97,0.97},
  xleftmargin=1em,
  xrightmargin=1em,
  aboveskip=0.8em,
  belowskip=0.8em,
  breaklines=true,
  showstringspaces=false,
  tabsize=2
}

\title{miniFlink: Declarative Stream Processing Semantics in OCaml}

\author{miniFlink Project}
\affiliation{%
  \institution{Independent}
  \country{Russia}}

\begin{abstract}
We present miniFlink, a functional implementation of Apache Flink semantics
in OCaml that demonstrates how declarative stream processing reduces
accidental complexity. The key insight is the \texttt{KEYED} type class:
by encoding the grouping key once in the type, operators \texttt{window},
\texttt{enrich}, and \texttt{dedup} require no explicit \texttt{\textasciitilde{}key:}
parameter, making the pipeline read as a specification rather than
an implementation. We implement event time, watermarks, tumbling/sliding
windows, stateful operators, exactly-once checkpointing, table and join
semantics, retractions for late data, and production layers
(dead letter queue, graceful shutdown, Prometheus metrics) in
approximately 1200 lines across 42 files. A parallel backend supports
both OCaml~4 (Thread + Mutex) and OCaml~5 (Domain + Atomic)
with measured speedup of $2.48\times$ on 4 workers (1 CPU).
\end{abstract}

\begin{document}
\maketitle

%%% ── 0. Declarativity ────────────────────────────────────────────────────
\section{Declarativity as a First Principle}

Traditional stream processing frameworks require the programmer to
explicitly state how to group events, how to serialize them, and how to
handle null values at every operator. Consider the same pipeline in
Kafka Streams (Java) and miniFlink (OCaml):

\begin{lstlisting}[caption={Kafka Streams (Java) — imperative},label={lst:kafka}]
KStream<String, Telemetry> stream =
    builder.stream("vehicles.telemetry");
KGroupedStream<String, Telemetry> grouped = stream
    .filter((k, v) -> v != null)
    .groupBy((k, v) -> v.getDeviceId(),
             Grouped.with(Serdes.String(), telemetrySerde));
KTable<Windowed<String>, Stats> stats = grouped
    .windowedBy(TimeWindows.ofSizeWithNoGrace(
                    Duration.ofSeconds(30)))
    .aggregate(Stats::empty,
               (k, v, acc) -> acc.add(v),
               Materialized.as("stats-store"));
stats.toStream()
    .flatMapValues(Rules::check)
    .filter((k, v) -> v != null)
    .to("fleet.alerts");
\end{lstlisting}

\begin{lstlisting}[caption={miniFlink (OCaml) — declarative},label={lst:mf}]
let pipeline source =
  source
  |> Event.with_watermarks   ~latency:(seconds 3)
  |> Pipe.enrich (module Telemetry)
       ~from:devices
       ~merge:(fun t dev -> { t with device = dev })
  |> Pipe.window (module Telemetry) (tumbling (seconds 30))
  |> Pipe.aggregate              Rules.compute
  |> Pipe.flat_map               (Rules.check Rules.fleet)
  |> Pipe.dedup (module Alert)
       ~rule:(fun a -> a.rule)
       ~cooldown:(minutes 5)
\end{lstlisting}

The differences are structural, not cosmetic:

\begin{itemize}
\item \textbf{No \texttt{key\_by}.} The \texttt{KEYED} type class encodes
  the grouping key in the type once; \texttt{window}, \texttt{enrich},
  and \texttt{dedup} extract it automatically.
\item \textbf{No type annotations} inside the \texttt{|>} chain.
  Hindley-Milner inference propagates types through all operators.
\item \textbf{No null checks.} \texttt{Option} makes absence explicit;
  \texttt{None} cannot be dereferenced.
\item \textbf{Format in one place.} Changing JSON to Protobuf means
  replacing one \texttt{Codec} value; the pipeline is untouched.
\end{itemize}

%%% ── 1. Architecture ─────────────────────────────────────────────────────
\section{Architecture}

miniFlink is organized into four orthogonal layers. Each layer depends
only on layers below it; the core never imports reliability or
parallelism modules.

\begin{table}[h]
\centering
\caption{Layer organization}
\begin{tabular}{lll}
\toprule
Layer & Modules & Responsibility \\
\midrule
Core & \texttt{stream, event, pipe, keyed, table, codec} &
  Immutable. No external dependencies. \\
Domain & \texttt{domain, rules} &
  \texttt{[@@deriving yojson]}. Pure functions. \\
Parallelism & \texttt{channel\_v4/v5, parallel\_v4/v5} &
  OCaml~4: Thread+Mutex. OCaml~5: Domain+Atomic. \\
Reliability & \texttt{dlq, shutdown, metrics, barrier, schema} &
  Each: \texttt{.mli} + \texttt{\_noop} + \texttt{\_default}. \\
\bottomrule
\end{tabular}
\end{table}

The \texttt{Makefile} detects \texttt{OCAML\_MAJOR} and copies
the appropriate \texttt{channel\_v\{4,5\}.ml} and
\texttt{parallel\_v\{4,5\}.ml} before compilation.

%%% ── 2. Core Abstractions ────────────────────────────────────────────────
\section{Core Abstractions}

\subsection{Pull Stream}

\begin{lstlisting}
type 'a t = unit -> 'a option
\end{lstlisting}

A stream is a function. The call stack \emph{is} the operator graph.
No callbacks, no threads, no allocation per event beyond the value
itself. \texttt{map}, \texttt{filter}, \texttt{flat\_map} compose
without intermediate buffers.

\subsection{Event}

\begin{lstlisting}
type 'a t =
  | Data      of 'a * timestamp   (* value + event time *)
  | Watermark of timestamp        (* lower bound on future ts *)
  | Retract   of 'a * timestamp   (* revocation of prior Data *)
\end{lstlisting}

All three constructors flow through the same stream. Operators
match exhaustively — adding a constructor forces all operators
to handle it at compile time.

\subsection{KEYED Type Class}

\begin{lstlisting}
module type S = sig
  type t
  val key : t -> string
end

(* Registered once per domain type *)
module Telemetry = Keyed.Make(struct
  type t = telemetry
  let key t = t.device_id
end)
\end{lstlisting}

Operators receive \texttt{(module Telemetry)} as a first-class module
and call \texttt{K.key} internally. The user never writes
\texttt{\textasciitilde{}key:(fun t -> t.device\_id)} again.

\subsection{Window}

The locally abstract type annotation breaks the occurs-check that
would otherwise create the recursive type
$\alpha = \mathtt{string} \times \alpha\ \mathtt{list}$:

\begin{lstlisting}
let window (type a) ?(latency=0)
    ~(key : a -> string) spec
    (upstream : a Event.t Stream.t)
    : (string * a list) Event.t Stream.t = ...
\end{lstlisting}

Window state is a \texttt{WMap.t} from \texttt{(key, start, stop)}
triples to event lists. Watermarks trigger partition of open windows
into \emph{closed} (fired) and \emph{open} sets; closed windows emit
their accumulated events downstream.

%%% ── 3. Event Time and Watermarks ───────────────────────────────────────
\section{Event Time and Watermarks}

Each event carries its origin timestamp independently of processing
time. \texttt{with\_watermarks} inserts \texttt{Watermark(max\_seen - latency)}
after each new maximum timestamp. The invariant:

\[
  \forall t < \mathtt{wm}:\ \text{no future } \mathtt{Data}(\_,\, t)
  \text{ will arrive}
\]

This allows deterministic window closure without wall-clock polling.
Watermarks are monotone by construction — \texttt{max\_seen} is a
\texttt{ref} that only grows.

%%% ── 4. Exactly-Once ─────────────────────────────────────────────────────
\section{Exactly-Once Checkpointing}

A checkpoint is a triple $(epoch, offset, state)$ stored in
\texttt{Checkpoint.Store}. The runner:

\begin{enumerate}
\item Reads $N$ events from source.
\item On failure: restores state from latest checkpoint,
  calls \texttt{source.seek(offset)}, resumes.
\item On success every $K$ events: atomically saves
  $(epoch{+}1,\ \mathtt{get\_offset}(),\ \mathtt{get\_state}())$.
\end{enumerate}

The key invariant is \emph{ordered commit}: flush sink before
writing checkpoint; write checkpoint before committing source offset.
Violated ordering creates either duplicates (checkpoint before flush)
or gaps (offset before checkpoint).

%%% ── 5. Table and Join ───────────────────────────────────────────────────
\section{Table and Join}

\begin{lstlisting}
type ('k,'v) table = 'k -> 'v option

let of_list pairs = Hashtbl.find_opt (Hashtbl.of_seq ...)
let of_stream ~key src = fun k -> pump src; Hashtbl.find_opt h k
\end{lstlisting}

A table is a function. \texttt{Snapshot} tables call \texttt{pump}
on each lookup to drain pending updates from an upstream stream.
\texttt{lookup\_join} is a one-liner:

\begin{lstlisting}
let lookup_join (module K) ~from ~merge upstream =
  emap (fun v -> merge v (from (K.key v))) upstream
\end{lstlisting}

\texttt{interval\_join} buffers both streams; events match when
keys are equal and $|t_A - t_B| \leq \mathtt{window\_ms}$. The
combined watermark is $\min(\mathtt{wm}_A,\, \mathtt{wm}_B)$.

%%% ── 6. Retractions ──────────────────────────────────────────────────────
\section{Retractions for Late Data}

When a late \texttt{Data} arrives after a window has fired, the
operator emits:
\[
  \mathtt{Retract}(v_\text{old},\, t) \;\to\; \mathtt{Data}(v_\text{new},\, t)
\]

Two aggregator strategies:

\begin{description}
\item[INVERTIBLE] $\mathtt{remove}(\mathtt{acc}, x)$ is defined.
  Late data: $O(1)$ update via \texttt{add}/\texttt{remove}.
  Examples: sum, count, average.
\item[RECOMPUTE] No inverse. Late data triggers full recomputation
  over the buffered event list: $O(n)$.
  Examples: max, min, percentile.
\end{description}

Both strategies produce identical external behaviour: the downstream
sees a \texttt{Retract}/\texttt{Data} pair and can perform an upsert.

%%% ── 7. Orthogonal Runtime ───────────────────────────────────────────────
\section{Orthogonal Runtime}

Each production capability has three files:
\texttt{foo.mli} (contract), \texttt{foo\_noop.ml} (zero-overhead stub),
\texttt{foo\_default.ml} (production implementation).
\texttt{runtime.ml} assembles them by \texttt{mode}:

\begin{lstlisting}
type mode = Noop | Log | Parallel | Prod

(* Noop: zero overhead, for tests *)
(* Log:  stderr metrics + SIGTERM handler *)
(* Prod: Prometheus HTTP :9090 + shutdown + DLQ log *)
\end{lstlisting}

\textbf{Dead Letter Queue.} \texttt{safe\_decode} wraps every codec
call; decode errors send the raw bytes to the DLQ with topic, error
message, and timestamp instead of crashing the pipeline.

\textbf{Graceful Shutdown.} A \texttt{guarded\_source} checks
\texttt{is\_requested()} before each poll. On SIGTERM: source returns
\texttt{None}, pipeline drains its buffer, sink flushes, process exits.

\textbf{Prometheus Metrics.} \texttt{metrics\_log} implements
counter, gauge, and histogram (7 buckets, $1\mu s$--$1s$) with
thread-safe \texttt{Mutex}-protected state and a TCP HTTP server
that serves the exposition text format on \texttt{GET /metrics}.

%%% ── 8. Parallelism ──────────────────────────────────────────────────────
\section{Parallelism}

Data parallelism via key sharding: \texttt{hash(key) mod N} routes
each event to a dedicated worker. Workers share no state — window
correctness follows from the invariant that all events for a given
key go to the same shard.

\begin{table}[h]
\centering
\caption{Benchmark: 1M events, 1000 devices, tumbling 30s, OCaml~4, 1 CPU}
\begin{tabular}{lrrr}
\toprule
Mode & ev/s & Time (s) & Speedup \\
\midrule
Sequential    & 209{,}293 & 4.78 & 1.00$\times$ \\
1 worker      & 151{,}570 & 6.60 & 0.72$\times$ \\
2 workers     & 405{,}775 & 2.46 & 1.94$\times$ \\
4 workers     & 519{,}463 & 1.93 & 2.48$\times$ \\
8 workers     & 599{,}988 & 1.67 & 2.87$\times$ \\
\bottomrule
\end{tabular}
\end{table}

The sub-linear speedup on 1 CPU is expected: OCaml~4's GIL prevents
true parallelism, but channel-mediated pipelining hides latency.
On OCaml~5 with \texttt{Domain}, channel overhead drops from
$\sim$106\,ns (Mutex) to $\sim$30\,ns (Atomic), and real
4-core parallelism is expected to yield $3.5$--$3.8\times$.

%%% ── 9. Testing ──────────────────────────────────────────────────────────
\section{Testing}

\textbf{Harness.} A deterministic test context \texttt{('a,'b) t}
feeds events via \texttt{push}/\texttt{push\_wm} and collects
output via \texttt{run}. No broker, no threads, no wall clock.

\textbf{Unit tests} (22): stream laws (map identity, filter
composition, flat\_map count), watermark monotonicity, window
correctness, dedup cooldown, enrich left-join semantics, rule
thresholds, full pipeline integration.

\textbf{Property tests} (7 × 100--500 cases via QCheck):

\begin{enumerate}
\item $\sum_w \mathtt{count}(w) = |\mathtt{input}|$
  — window preserves all events.
\item $\mathtt{dedup}(\mathtt{dedup}(s)) = \mathtt{dedup}(s)$
  — dedup is idempotent.
\item Watermarks are monotone for any event order.
\item \texttt{enrich} preserves event count (left join).
\item $\mathtt{critical} \Leftrightarrow \mathtt{speed} > 1.3 \times \mathtt{limit}$
  — for all \texttt{float} pairs.
\item \texttt{map} preserves event count.
\item $|\mathtt{filter}(s)| \leq |s|$.
\end{enumerate}

%%% ── 10. Conclusion ──────────────────────────────────────────────────────
\section{Conclusion}

miniFlink demonstrates that the core semantics of Apache Flink —
event time, watermarks, windowing, stateful operators, exactly-once,
table/join, retractions, and production reliability layers — fit
in $\sim$1200 lines of OCaml. Two design decisions proved most
consequential:

\begin{enumerate}
\item The \texttt{KEYED} type class eliminates repeated key
  extraction, making the pipeline read as a specification.
\item The \texttt{(type a)} locally abstract type annotation in
  \texttt{Pipe.window} avoids the recursive type $\alpha = \mathtt{string}
  \times \alpha\ \mathtt{list}$ that otherwise requires \texttt{Obj.magic}.
\end{enumerate}

The OCaml module system's ability to accept modules as runtime values
(first-class modules) is what makes both design decisions expressible
without metaprogramming or code generation.

\end{document}
